\documentclass[a4paper,12pt]{report}
\usepackage[utf8]{inputenc}
\usepackage[T1]{fontenc}
\usepackage[french]{babel}
\usepackage{graphicx}
\usepackage{amsfonts}
\usepackage{float}
\usepackage[export]{adjustbox}
\usepackage{tikz}
\usepackage{algorithm}
\usepackage{algpseudocode}
\usepackage{subfig}
\usepackage[hidelinks]{hyperref}
\usepackage{listings}
\usepackage{xcolor}

\colorlet{punct}{red!60!black}
\definecolor{background}{HTML}{EEEEEE}
\definecolor{delim}{RGB}{20,105,176}
\colorlet{numb}{magenta!60!black}

\lstdefinelanguage{json}{
    basicstyle=\normalfont\ttfamily,
    numbers=left,
    numberstyle=\scriptsize,
    stepnumber=1,
    numbersep=8pt,
    showstringspaces=false,
    breaklines=true,
    frame=lines,
    backgroundcolor=\color{background},
}


\lstset{frame=tb,
  language=json,
  aboveskip=3mm,
  belowskip=3mm,
  showstringspaces=false,
  columns=flexible,
  basicstyle={\small\ttfamily},
  numbers=none,
  breaklines=true,
  breakatwhitespace=true,
  tabsize=3
}

\date{}
\title{Rapport AWS \\ Groupe 8 : Projet TODO LIST}

\author{Niels Merceron \\ Pierre Vermeulen \\ Alexis Guigal \\ Manel Azgag \\ \\  \includegraphics[scale=0.20]{logo-UVSQ-2020-RVB.png}}


\begin{document}
\maketitle

\newpage
\tableofcontents
\chapter{Introduction}

Nous sommes un groupe avec peu de notions de web, donc pour nous familiariser avec les langages et autre technologie nous avons fait le choix de faire un site simple et plus particulièrement un site de to do list. Il y a beaucoup de ressource en ligne nous décrivant comment faire une to do list et donc nous aiderons tout au cours du développement de notre site pour savoir quel service implémenter de notre site ou non, si c'est à notre portée.
\\Pour autant, faire ce site relèvera de quelques challenges de conception vu que nous ne sommes pas familiers avec beaucoup de technologie pour développer des sites web.
On peut noter par exemple la gestion du back peut être vite contraignante, ou l'aspect sécurité qui doit être fait de manière la plus propre possible pour ne laisser aucune faille.


\chapter{Technologie utilisée}
\section{Front end}

\subsection{Svelte}
C'est un framwork simple et proche syntaxiquement du javascript standard. Il est conçu pour compiler le code au moment de la compilation plutôt qu'au moment de l'éxécution, il offre donc de très bonnes performances. Le framework est également bien documenté avec une documentation claire et précise.

\subsection{Tailwind avec DaisyUI}
C'est un framwork qui est bien plus simple que les autres framework CSS.Sa documentation , sa simplicité d'utilisation et de la modernité des design présent fait de ce framework le meilleur choix pour notre application.Il est également très facile à intégrer avec n'importe quel type de stack.

\section{Back end}

\subsection{Nodejs/expressjs}
C'est un environnement d'éxécution pour JavaScript construit sur le moteur Javascript de Chrome. Il est extrêmement populaire et il est utilisé par de très nombreuses entreprises. Il est également très simple à prendre en main.

\subsection{JWT}

Le jsonwebtoken est un standard définit dans la RFC 7519. 
La technologie de jsonwebtoken est publique (c'est une norme ouverte) qui a pour but d'avoir un format compact et autonome pour assurer des communications sécurisées entre deux ou plusieurs éléments.
Cette technologie existe depuis 2015.\\
Un JWT est composé de la manière suivante:

\begin{itemize}
\item un header contenant l'algorithme de chiffrement, le type de token (JWT)
\item un payload (la charge utile) qui contient toutes les informations transmises à l'application
\item Une signature qui est générée en fonction du header ,du payload et d'une clef secrète connu que part l'application.\\
\end{itemize}

On peut l'utiliser dans à peut près tous les cas possibles car il est assez "fléxible" de conception. 
Pour notre projet nous utiliserons le JWT lié a des cookies pour éviter que le contenue du JWT soit accessible par l'utilisateur.

\section{Sécurité}
Pour la sécurité, nous avons utilisé plusieurs bibliothèques natives de nodejs mais aussi certaines que nous avons vu pendant le cours ou décourvertes durant les recherches pour ce projet.

\subsection{Bcrypt et crypto}
 On utilse Bcrypt pour générer des hashé de mot de passe qui est basé sur l'algorithme du "blowfish cipher". Pour éviter des attaques par force brute nous ajoutons une entré aléatoire a Bcrypt. On utilise Bcrypt principalement dans le back.\\
 
 Crypto est une bibliothèque que nous avons voulu utiliser car elle est assez versatille. Elle peut aussi bien nous être utilse pour générer des Hash,que de signer opu faire des vérifications sur des données cryptées. Nous l'avons utilisé dans le Back et le Front du site pour les connexions et stockage des mots de passe.
 
\subsection{Helmet}
	Helmet est une bibliothèque native de nodejs qui nous permet de sécurisé avec des headers les communications entre le back est le front. Ces headers nou sprotège par exemple des attaque de "cross-site-Scripting"

\subsection{Protocole TLS}

\section{SGBD}
\subsection{SQL ou NoSQL}
Nous avions donc opter pour un SGBD NoSQL malgré le fait qu'on ai plus pratiqué le SGBD relationnel. Cependant le NoSQL est plus adapté et plus le projet évolue, plus le NoSQL deviennait une évidence.
\subsection{Mongodb}
MongoDB est un système de gestion de bases de données NoSQL, qui utilise un format de stockage de données basé sur JSON. Contrairement aux bases de données relationnelles, MongoDB n'utilise pas de tables et de schémas fixes, mais stocke les données dans des collections flexibles qui peuvent être modifiées sans avoir à définir un schéma préalablement.\\
Cette approche permet une grande flexibilité pour gérer des données complexes et des schémas évolutifs, ce qui est particulièrement utile pour les applications modernes basées sur le cloud. De plus, MongoDB offre des performances élevées, une évolutivité horizontale facile et une grande disponibilité, ce qui en fait une solution de choix pour les entreprises qui cherchent à développer des applications à grande échelle.
En raison de ses avantages, MongoDB est très utilisé et dispose d'une documentation de qualité, ce qui facilite son adoption et son utilisation.


\chapter{Déroulement du projet}

\section{Back}
\subsection{Création des routes}
\subsubsection{- Signup - POST}
Exemple de requête:
\begin{lstlisting}
{
	"username": "bonjour",
	"email": "test@fafa.com",
	"password": "drffffffff"
}
\end{lstlisting}

Exemple de réponse:
\begin{lstlisting}
{
	"token": "eyJhbGciOiJIUzI1Ni5cCI6...2v0RWSOkv4"
}
\end{lstlisting}

\subsubsection{- Login - POST}
Exemple de requête:
\begin{lstlisting}
{
  	"email": "test@fafa.com",
  	"password": "drffffffff"
}
\end{lstlisting}

Exemple de réponse:
\begin{lstlisting}
{
	"token": "eyJhbGciOiJIUzI1R5cCI6...2v0RWSOkv4"
}
\end{lstlisting}
\subsubsection{- Createtodo - POST}
Exemple de requête:
\begin{lstlisting}
{
	"title": "tesvqhtret1",
	"description": "gfzghfgrfeafezrgezg"
}
\end{lstlisting}

Exemple de réponse:
\begin{lstlisting}
	{"title":"tesvqhtret1",
	description":"gfzghfgrfeafezrgezg",
	"completed":false,
	"createdAt":"2023-04-28T23:57:41.194Z",
	_id":"644c5ef0c7ff2ca10c44521c","__v":0}
\end{lstlisting}

\subsubsection{- Gettodo - GET}
Exemple de requête:


Exemple de réponse:
\begin{lstlisting}
[
	{
	"_id": "644c5647cd148f99d95d98a7",
	"title": "test1",
	description": "gfzghezrgezg",
	"completed": false,
	"createdAt": "2023-04-28T23:26:51.777Z",
	"__v": 0
	},
	{
	"_id": "644c5ef0c7ff2ca10c44521c",
	"title": "YES",
	"description": "mod",
	"completed": false,
	"createdAt": "2023-04-28T23:57:41.194Z",
	"__v": 0
	}
]
\end{lstlisting}

\subsection{- Del - DEL}
Exemple de requête:

Exemple de réponse:
\begin{lstlisting}
{
	"message": "Todo Deleted"
}
\end{lstlisting}

\subsection{- Modify - PUT}
Exemple de requête:

Exemple de réponse:
\begin{lstlisting}
{
	"_id": "644c5ef0c7ff2ca10c44521c",
	"title": "YES",
	"description": "mod",
	"completed": false,
	"createdAt": "2023-04-28T23:57:41.194Z",
	"__v": 0
}
\end{lstlisting}

\chapter{Objectif accompli durant le projet}

Dans cette section, on vous présente l'état final de notre projet.



\chapter{Conclusion}
Aprés 2 mois et demi de travail en role alternés, nous sommes fiers de vous présenter notre site de todolist.
Ce projet nous a permis de développer de nouvelles compétences en matière de développement web,
d'expérimenter de nouvelles technologies et de relever des défis d'implémentation et d'organisation.
Nous sommes heureux d'avoir pu mettre en pratique les concepts que nous avons appris et de voir 
notre travail prendre vie sous la forme d'un outil utile et pratique pour les utilisateurs. 
Nous espérons que notre site de todolist sera une aide précieuse pour la gestion de tâches quotidiennes et nous
sommes impatients de continuer à améliorer l'expérience utilisateur dans les futures versions.
\end{document}
